% --------------------------------------------------
% Chapter 05
% MLOps vs DevOps vs DataOps
% --------------------------------------------------

\label{chap:mlops-devops-dataops}

\section{Learning Objectives}

After completing this chapter, learners will be able to:

\begin{itemize}
    \item Define DevOps, DataOps, and MLOps.
    \item Understand the objectives of each discipline.
    \item Identify similarities and differences.
    \item Explain how the three disciplines complement each other.
    \item Design collaborative workflows across teams.
\end{itemize}

\section{Introduction}

Modern organizations rely on multiple operational disciplines to deliver reliable digital products and data-driven services.

Among the most important disciplines are:

\begin{itemize}
    \item DevOps
    \item DataOps
    \item MLOps
\end{itemize}

Although these disciplines share common principles, they address different challenges and focus on different assets.

Understanding their respective roles is essential for building enterprise-grade machine learning platforms.

\section{DevOps}

\subsection{Definition}

DevOps is a set of practices that combines software development and IT operations.

Its primary objective is to improve the speed, reliability, and quality of software delivery.

DevOps emerged to eliminate the traditional separation between development teams and operations teams.

\subsection{Core Objectives}

DevOps seeks to:

\begin{itemize}
    \item Accelerate software delivery
    \item Improve deployment reliability
    \item Reduce operational incidents
    \item Increase collaboration
    \item Automate repetitive tasks
\end{itemize}

\subsection{Typical DevOps Activities}

Examples include:

\begin{itemize}
    \item Source code management
    \item Automated testing
    \item Continuous Integration
    \item Continuous Delivery
    \item Infrastructure management
    \item Application monitoring
\end{itemize}

\subsection{Primary Assets}

DevOps primarily manages:

\begin{itemize}
    \item Source code
    \item Applications
    \item Infrastructure
    \item Deployment environments
\end{itemize}

\subsection{Typical Tooling}

Examples include:

\begin{itemize}
    \item Git
    \item GitHub
    \item Jenkins
    \item GitHub Actions
    \item Docker
    \item Kubernetes
\end{itemize}

\section{DataOps}

\subsection{Definition}

DataOps applies DevOps-inspired principles to data management and analytics workflows.

Its objective is to improve the reliability, quality, and efficiency of data pipelines.

DataOps recognizes that machine learning systems are only as good as the data they consume.

\subsection{Core Objectives}

DataOps aims to:

\begin{itemize}
    \item Improve data quality
    \item Automate data pipelines
    \item Increase data reliability
    \item Improve collaboration between data teams
    \item Reduce data delivery times
\end{itemize}

\subsection{Typical DataOps Activities}

Examples include:

\begin{itemize}
    \item Data ingestion
    \item Data transformation
    \item Data validation
    \item Data quality monitoring
    \item Metadata management
    \item Data lineage tracking
\end{itemize}

\subsection{Primary Assets}

DataOps primarily manages:

\begin{itemize}
    \item Datasets
    \item Data pipelines
    \item Metadata
    \item Data quality controls
\end{itemize}

\subsection{Typical DataOps Tooling}

Examples include:

\begin{itemize}
    \item Apache Airflow
    \item Azure Data Factory
    \item dbt
    \item Databricks
    \item Spark
    \item Great Expectations
\end{itemize}

\section{MLOps}

\subsection{Definition}

MLOps extends DevOps and DataOps concepts to machine learning systems.

Its objective is to operationalize machine learning models throughout their lifecycle.

Unlike traditional software systems, machine learning systems introduce additional challenges related to:

\begin{itemize}
    \item Data dependency
    \item Model degradation
    \item Experiment management
    \item Retraining requirements
    \item Governance obligations
\end{itemize}

\subsection{Core Objectives}

MLOps aims to:

\begin{itemize}
    \item Improve reproducibility
    \item Automate model development
    \item Standardize deployment
    \item Monitor model behavior
    \item Support governance and compliance
\end{itemize}

\subsection{Typical MLOps Activities}

Examples include:

\begin{itemize}
    \item Experiment tracking
    \item Model versioning
    \item Training pipelines
    \item Model deployment
    \item Drift monitoring
    \item Retraining automation
\end{itemize}

\subsection{Primary Assets}

MLOps primarily manages:

\begin{itemize}
    \item Models
    \item Features
    \item Experiments
    \item Training pipelines
    \item Model registries
\end{itemize}

\subsection{Typical MLOps Tooling}

Examples include:

\begin{itemize}
    \item MLflow
    \item Azure Machine Learning
    \item Kubeflow
    \item Feast
    \item Weights \& Biases
\end{itemize}

\section{Comparing DevOps, DataOps, and MLOps}

\subsection{High-Level Comparison}

\begin{table}[h]
\centering
\begin{tabular}{p{3cm}p{3cm}p{3cm}p{3cm}}
\toprule
Aspect & DevOps & DataOps & MLOps \\
\midrule
Primary Asset & Code & Data & Models \\
Main Goal & Software Delivery & Data Reliability & Model Operations \\
Versioned Objects & Code & Datasets & Models and Data \\
Monitoring Focus & Applications & Pipelines & Models \\
Deployment Target & Applications & Data Pipelines & ML Services \\
\bottomrule
\end{tabular}
\caption{High-level comparison of DevOps, DataOps, and MLOps}
\end{table}

\subsection{Common Principles}

Although they focus on different assets, all three disciplines share several principles:

\begin{itemize}
    \item Automation
    \item Version Control
    \item Monitoring
    \item Collaboration
    \item Continuous Improvement
\end{itemize}

\subsection{Key Differences}

The main differences arise from the nature of the managed assets.

Software applications are deterministic.

Datasets evolve continuously.

Machine learning models depend on both code and data.

This additional complexity explains why MLOps is often considered a superset of both DevOps and DataOps concepts.

\section{The Additional Complexity of MLOps}

\subsection{Model Drift}

Traditional applications do not typically degrade over time.

Machine learning models may experience:

\begin{itemize}
    \item Data Drift
    \item Concept Drift
    \item Performance Degradation
\end{itemize}

This requires continuous monitoring.

\subsection{Experiment Management}

Software engineering rarely involves hundreds of competing versions of the same application.

Machine learning development frequently involves:

\begin{itemize}
    \item Multiple algorithms
    \item Multiple datasets
    \item Multiple feature sets
    \item Multiple hyperparameter configurations
\end{itemize}

Experiment tracking becomes essential.

\subsection{Retraining Requirements}

Applications are usually updated when requirements change.

Machine learning models often require retraining even when requirements remain unchanged.

This introduces additional operational complexity.

\section{Collaboration Model}

\subsection{Why Collaboration Matters}

Modern machine learning systems require expertise from multiple domains.

No single team can successfully deliver and maintain production-grade machine learning systems independently.

Successful organizations establish collaboration mechanisms across software, data, and machine learning teams.

\subsection{Typical Responsibilities}

\begin{table}[h]
\centering
\begin{tabular}{p{4cm}p{8cm}}
\toprule
Role & Primary Responsibilities \\
\midrule
Data Engineer & Data pipelines, ingestion, transformation \\
Data Scientist & Model development and experimentation \\
ML Engineer & Model operationalization and deployment \\
DevOps Engineer & Infrastructure and CI/CD \\
Product Owner & Business requirements and prioritization \\
Model Validator & Independent model review \\
Risk Manager & Governance and compliance \\
\bottomrule
\end{tabular}
\caption{Typical responsibilities in machine learning organizations}
\end{table}

\subsection{Integrated Workflow}

A simplified workflow may be represented as:

\begin{center}
Business Requirements
$\rightarrow$
Data Engineering
$\rightarrow$
Model Development
$\rightarrow$
Validation
$\rightarrow$
Deployment
$\rightarrow$
Monitoring
\end{center}

Each stage requires collaboration across multiple functions.

\section{Organizational Structures}

\subsection{Traditional Structure}

Historically, organizations separated:

\begin{itemize}
    \item Software Development
    \item Data Engineering
    \item Analytics
    \item Operations
\end{itemize}

This often resulted in communication gaps and delayed deployments.

\subsection{Cross-Functional Teams}

Modern organizations increasingly adopt cross-functional teams composed of:

\begin{itemize}
    \item Data Scientists
    \item Data Engineers
    \item ML Engineers
    \item DevOps Engineers
    \item Business Stakeholders
\end{itemize}

Benefits include:

\begin{itemize}
    \item Faster delivery
    \item Improved collaboration
    \item Better ownership
    \item Reduced handoffs
\end{itemize}

\section{Credit Scoring Case Study}

\subsection{Scenario}

A bank develops a new credit scoring platform.

The project involves multiple stakeholders.

\subsection{DataOps Contributions}

The DataOps team is responsible for:

\begin{itemize}
    \item Collecting customer data
    \item Integrating bureau information
    \item Building data pipelines
    \item Ensuring data quality
\end{itemize}

\subsection{MLOps Contributions}

The MLOps team is responsible for:

\begin{itemize}
    \item Experiment tracking
    \item Model training pipelines
    \item Model registry management
    \item Deployment automation
    \item Drift monitoring
\end{itemize}

\subsection{DevOps Contributions}

The DevOps team is responsible for:

\begin{itemize}
    \item Infrastructure provisioning
    \item Container orchestration
    \item CI/CD pipelines
    \item Platform monitoring
    \item Security controls
\end{itemize}

\subsection{Outcome}

The success of the platform depends on effective collaboration across all three disciplines.

None of the disciplines can independently deliver a complete production-grade solution.

\section{Practical Guidance}

Organizations beginning their MLOps journey should:

\begin{enumerate}
    \item Establish source control practices.
    \item Automate data pipelines.
    \item Implement experiment tracking.
    \item Introduce CI/CD pipelines.
    \item Deploy monitoring solutions.
    \item Strengthen governance processes.
\end{enumerate}

Maturity typically increases gradually rather than through a single transformation initiative.

\section{Summary}

DevOps, DataOps, and MLOps are complementary disciplines.

Each focuses on different assets:

\begin{itemize}
    \item DevOps manages software and infrastructure.
    \item DataOps manages data and pipelines.
    \item MLOps manages models and machine learning workflows.
\end{itemize}

All three disciplines share common principles including automation, monitoring, collaboration, and continuous improvement.

Successful machine learning organizations integrate these disciplines into a unified operating model capable of delivering reliable, scalable, and governable AI systems.

The next chapter introduces the MLOps tooling landscape and presents the major technologies used throughout the machine learning lifecycle.